\begin{ouverture}[college1]{2}{Atomes, noyaux et ions}{Partie 1 · Organisation et transformations de la matière}{Comment une entité peut-elle devenir chargée ?}
\objectif{Décrire les constituants de l’atome.}
\objectif{Lire Z et un inventaire de nucléons.}
\objectif{Interpréter une formule moléculaire et une charge ionique.}
\prerequis{Symboles chimiques et nombres entiers.}
\prerequis{Charges positives et négatives.}
\end{ouverture}
\pcPage{Le cours}{Atomes, noyaux et ions}
\phantomsection\label{ess-2}
\essSection{1}{Un noyau et des électrons}
Un atome comporte un noyau central et un nuage d’électrons. Le noyau contient des protons, de charge positive, et généralement des neutrons, sans charge électrique. Les électrons portent une charge négative opposée à celle du proton. Un atome neutre possède autant de protons que d’électrons. Sa masse est presque entièrement concentrée dans le noyau. Le dessin de collège ne représente pas des orbites réelles.\par
\pcFigure{figures/b02-atome.pdf}{Modèle schématique : noyau très agrandi ; les électrons sont figurés par des points.}
\essSection{2}{Le numéro atomique}
Le numéro atomique Z est le nombre de protons et identifie l’élément. Si A est le nombre total de nucléons, le nombre de neutrons vaut $A-Z$. Cette notation est exploitée lorsqu’elle est donnée. Pour $Z=6$ et $A=12$, le noyau contient six protons et six neutrons ; l’atome neutre possède six électrons. Un tableau périodique permet de retrouver nom, symbole et Z.\par
\clearpage
\essSection{3}{Former un ion}
Un ion est une entité chargée ayant perdu ou gagné un ou plusieurs électrons. Une perte d’électrons donne une charge positive ; un gain donne une charge négative. Lors d’une ionisation chimique usuelle, le noyau conserve son nombre de protons : l’élément ne change pas. $\ce{Na+}$ a perdu un électron ; $\ce{Cl-}$ en a gagné un ; $\ce{Cu^{2+}}$ en a perdu deux.\par
\pcFigure{figures/b02-ions.pdf}{Les nombres d’électrons indiqués concernent l’entité entière ; aucun proton n’est déplacé.}
\essSection{4}{Molécules, indices et coefficients}
Une molécule est un assemblage neutre d’atomes liés. $\ce{H2O}$ contient deux H et un O ; $\ce{CO2}$ un C et deux O ; $\ce{CH4}$ un C et quatre H ; $\ce{N2O}$ deux N et un O. Dans $3\,\ce{CO2}$, le coefficient multiplie toute la formule : trois C et six O. Un indice change l’espèce ; un coefficient change sa quantité d’entités. Un ion peut aussi être polyatomique.\par
\begin{pcEncadre}{method}{Un exemple expliqué}Le noyau d’un ion aluminium contient 13 protons et 14 neutrons ; l’entité possède 10 électrons. $Z=13$, $A=27$. La charge totale correspond à $13-10=3$ charges positives : $\ce{Al^{3+}}$. L’atome neutre d’aluminium possède 13 électrons.\end{pcEncadre}
\pcPage{Expérimenter et retenir}{Atomes, noyaux et ions}
\phantomsection\label{activite-2}
\begin{pcEncadre}{method}{Investigation en classe}À partir de cartes « noyau » et de jetons électrons, construire un atome neutre, puis un ion du même élément. Justifier chaque inventaire. Utiliser ensuite un tableau périodique pour vérifier les symboles. Séparer explicitement les limites du modèle graphique et les conclusions sur les charges.\end{pcEncadre}
\pcSousTitre{Ce que je dois savoir expliquer}\begin{questions}
\item Définir Z et distinguer proton, neutron et électron.
\item Retrouver les électrons d’un ion simple avec Z et sa charge.
\item Compter les atomes sans confondre indice et coefficient.
\end{questions}
\begin{pcEncadre}{attention}{Une erreur à éviter}Un ion positif n’a pas gagné des protons : dans les transformations chimiques étudiées, ce sont les électrons qui peuvent être échangés.\end{pcEncadre}
\pcRepere{cor-activite-2}{Pistes pour l’investigation}\pcGrille
\clearpage\markright{Je démarre}
\rubrique{Je démarre}
\par\noindent\begin{minipage}{\linewidth}
\begin{exercice}[1]{Constituants}
\begin{questions}
\item Où se trouvent protons et neutrons ?
\item Quelle particule porte une charge négative ?
\end{questions}
\end{exercice}
\end{minipage}
\par\noindent\begin{minipage}{\linewidth}
\begin{exercice}[1]{Neutralité}
\begin{questions}
\item Un atome possède 8 protons. Combien d’électrons ?
\item Quel est son numéro atomique ?
\end{questions}
\end{exercice}
\end{minipage}
\par\noindent\begin{minipage}{\linewidth}
\begin{exercice}[1]{Un noyau}
\begin{questions}
\item Un noyau contient 11 protons et 12 neutrons. Donner Z.
\item Donner A.
\end{questions}
\end{exercice}
\end{minipage}
\par\noindent\begin{minipage}{\linewidth}
\begin{exercice}[1]{L’eau}
\begin{questions}
\item Compter H et O dans une molécule d’eau.
\item Compter ces atomes dans trois molécules.
\end{questions}
\end{exercice}
\end{minipage}
\par\noindent\begin{minipage}{\linewidth}
\begin{exercice}[1]{Le signe}
\begin{questions}
\item L’ion $\ce{Cl-}$ a-t-il gagné ou perdu un électron ?
\item Même question pour $\ce{Na+}$.
\end{questions}
\end{exercice}
\end{minipage}
\clearpage\markright{Je m’entraîne}
\rubrique{Je m’entraîne}
\par\noindent\begin{minipage}{\linewidth}
\begin{exercice}[2]{Un ion cuivre}
\begin{questions}
\item Cu a $Z=29$. Combien d’électrons dans $\ce{Cu^{2+}}$ ?
\item Combien de protons dans cet ion ?
\end{questions}
\end{exercice}
\end{minipage}
\par\noindent\begin{minipage}{\linewidth}
\begin{exercice}[2]{Un ion chlorure}
\begin{questions}
\item Cl a $Z=17$. Compter les électrons de $\ce{Cl-}$.
\item Cet ion est-il une molécule de dichlore $\ce{Cl2}$ ?
\end{questions}
\end{exercice}
\end{minipage}
\par\noindent\begin{minipage}{\linewidth}
\begin{exercice}[2]{Inventaire complet}
\begin{questions}
\item Un atome a Z = 13 et A = 27. Compter neutrons et électrons.
\item Il perd trois électrons. Donner la charge et le symbole, sachant qu’il s’agit d’aluminium.
\end{questions}
\end{exercice}
\end{minipage}
\par\noindent\begin{minipage}{\linewidth}
\begin{exercice}[2]{Indices et coefficient}
\begin{questions}
\item Compter les atomes dans $2\,\ce{CH4}$.
\item Compter les atomes dans $3\,\ce{N2O}$.
\end{questions}
\end{exercice}
\end{minipage}
\par\noindent\begin{minipage}{\linewidth}
\begin{exercice}[2]{Trois écritures}
\begin{questions}
\item Distinguer O, $\ce{O2}$ et $\ce{O^{2-}}$.
\item Les trois appartiennent-ils au même élément ?
\end{questions}
\end{exercice}
\end{minipage}
\clearpage\markright{Je raisonne}
\rubrique{Je raisonne}
\par\noindent\begin{minipage}{\linewidth}
\begin{exercice}[3]{Retrouver un élément}
\begin{questions}
\item Une entité possède 12 protons et 10 électrons. Trouver sa charge.
\item Le tableau donne Mg pour Z = 12. Écrire l’ion.
\end{questions}
\end{exercice}
\end{minipage}
\par\noindent\begin{minipage}{\linewidth}
\begin{exercice}[3]{Une fausse justification}
\begin{questions}
\item « Le sodium devient chlorure parce qu’il gagne un électron. » Corriger cette phrase.
\item Quel nombre devrait changer pour changer d’élément ?
\end{questions}
\end{exercice}
\end{minipage}
\par\noindent\begin{minipage}{\linewidth}
\begin{exercice}[3]{Un changement d’échelle}
\begin{questions}
\item On prend un diamètre atomique de $10^{-10}\unite m$ et un diamètre nucléaire de $10^{-15}\unite m$. Calculer le rapport.
\item Si l’atome est représenté par 10 m, quel diamètre aurait le noyau ?
\end{questions}
\end{exercice}
\end{minipage}
\par\noindent\begin{minipage}{\linewidth}
\begin{exercice}[3]{Une solution électriquement neutre}
\begin{questions}
\item On représente deux ions $\ce{Na+}$ et deux $\ce{Cl-}$. Quelle charge totale ?
\item Avec deux $\ce{Ca^{2+}}$, combien de $\ce{Cl-}$ faut-il pour compenser ?
\end{questions}
\end{exercice}
\end{minipage}
\par\noindent\begin{minipage}{\linewidth}
\begin{exercice}[3]{Deux noyaux du même élément}
\begin{questions}
\item Deux noyaux ont Z = 6, mais A = 12 et A = 14. Comparer les nombres de neutrons.
\item Pourquoi sont-ils du même élément malgré cette différence ?
\end{questions}
\end{exercice}
\end{minipage}
\pcPage{Défi du chapitre}{Atomes, noyaux et ions}
\phantomsection\label{defi-ex-2}
\pcDefiTitre{2}{L’entité inconnue}
\begin{questions}
\item Un noyau contient 26 protons et 30 neutrons ; l’entité possède 23 électrons. Donner Z et A.
\item Le tableau donne Fe pour Z = 26. Écrire le symbole de l’ion.
\item Expliquer en trois phrases pourquoi il ne s’agit ni d’un nouvel élément ni d’une molécule.
\end{questions}
\pcGrille
